\section{数据来源与参考文献}

\begin{plainbox}
本系统所用数据均来自\textbf{公开学术数据源},算法与模型均为公开方法或开源权重,知识产权清晰、可追溯、可复现。
\end{plainbox}

\subsection{数据来源}
\begin{enumerate}[label=D\arabic*.]
\item OpenAlex —— 开放学术知识图谱与元数据 API。\texttt{https://openalex.org}
\item PubMed / PubMed Central (PMC) —— 美国国立医学图书馆生物医学文献库。\texttt{https://pubmed.ncbi.nlm.nih.gov}
\item Crossref —— DOI 注册与文献元数据。\texttt{https://www.crossref.org}
\item DOAJ —— 开放获取期刊目录。\texttt{https://doaj.org}
\item arXiv —— 计算机/物理/数学等领域预印本库(含全文 PDF)。\texttt{https://arxiv.org}
\end{enumerate}

\subsection{方法与模型来源}
\begin{enumerate}[label=R\arabic*.]
\item Mann, H. B. (1945). Nonparametric tests against trend. \textit{Econometrica}, 13(3), 245--259.
\item Sen, P. K. (1968). Estimates of the regression coefficient based on Kendall's tau. \textit{JASA}, 63(324), 1379--1389.
\item Blondel, V. D., et al. (2008). Fast unfolding of communities in large networks (Louvain). \textit{J. Stat. Mech.}
\item Page, L., Brin, S., et al. (1999). The PageRank citation ranking: Bringing order to the web.
\item Blei, D., Ng, A., Jordan, M. (2003). Latent Dirichlet Allocation (LDA). \textit{JMLR}, 3, 993--1022.
\item Lee, D. D., Seung, H. S. (1999). Learning the parts of objects by non-negative matrix factorization (NMF). \textit{Nature}, 401, 788--791.
\item Cormack, G. V., Clarke, C. L. A., Büttcher, S. (2009). Reciprocal rank fusion (RRF). \textit{SIGIR}.
\item Robertson, S., Zaragoza, H. (2009). The probabilistic relevance framework: BM25 and beyond.
\item Wang, L., et al. (2024). Multilingual E5 text embeddings(\texttt{multilingual-e5-base},检索与跨语言嵌入)。
\item Chen, J., et al. (2024). BGE M3-Embedding / \texttt{bge-reranker-v2-m3}(跨语言重排)。
\item Qwen Team (2024). Qwen2.5 Technical Report(本地生成层 \texttt{Qwen2.5-7B-Instruct})。
\end{enumerate}

\begin{plainbox}
\textbf{合规声明}:数据用于科研分析与方法验证,遵循各数据源的开放获取与使用条款;系统不存储或再分发受版权保护的论文全文,仅保留用于检索的元数据、摘要与开放获取正文片段。
\end{plainbox}
